\section{Results}
\subsection{Recovery changes with requested duration}
The released severity~2 recovery shows a large duration interaction. CLAP gain rises from $R=+0.085$ \ci{+0.066}{+0.105} at $3.84$\,s to $+0.244$ \ci{+0.215}{+0.273} at $10.24$\,s, giving $J=+0.159$ \ci{+0.131}{+0.187}. The intermediate points in Table~\ref{tab:evidence} show a gradual response rather than an endpoint anomaly. The scale also changes materially. Recovery closes $44\%$ of the gap from P to dense at $3.84$\,s and $82\%$ at $10.24$\,s.

Human-CLAP, KL to matched real references and PANNs top-10 capture all support larger recovery as duration increases, but they do not imply continued growth at every step. Human-CLAP is essentially saturated by $7.68$\,s, changing only from $0.371$ to $0.375$ recovery at $10.24$\,s. The published DDIM $200$, guidance $3.5$ recipe still reproduces the endpoint interaction with $J=+0.184$ \ci{+0.126}{+0.243}, and the registered severity~2 conclusions survive Holm correction. These checks reuse the AudioCaps prompt draw, so they establish metric and sampler robustness rather than an independent prompt replication.

Extending generation beyond the native duration does not reveal a sharp peak. On the matched $96$-prompt subset, recovery changes by $+0.021$ \ci{-0.023}{+0.067} from $10.24$ to $15.36$\,s. The interval is too wide to establish a plateau, but it provides no clear evidence of another increase comparable with the earlier sweep steps.

\subsection{Training duration does not determine the favorable evaluation duration}
The matched $20{,}000$-step intervention supplies the direct test. The checkpoint trained at $3.84$\,s has $J_{3.84}=+0.065$ \ci{+0.044}{+0.087}; the checkpoint trained at $10.24$\,s has the smaller interaction $J_{10.24}=+0.031$ \ci{+0.010}{+0.052}. Their paired difference is $\Delta J=-0.035$ \ci{-0.064}{-0.004}. Training-duration specialization predicts the opposite sign because moving training to $10.24$\,s should increase the long-minus-short advantage. At this matched budget, the favorable evaluation duration therefore does not follow the duration used for fine-tuning.

The dense $2\times2$ gives a clean negative result for this particular training recipe. Raw and EMA dense baselines differ by only $+0.001$ \ci{-0.017}{+0.020} at $3.84$\,s and $-0.017$ \ci{-0.044}{+0.010} at $10.24$\,s, whereas every $20{,}000$-step dense gain is strongly negative, between $-0.220$ and $-0.168$ with upper $95\%$ bounds below $-0.144$. Its training-duration contrast is nevertheless unresolved, $\Delta J_{\rm dense}=-0.005$ \ci{-0.033}{+0.025}. The quantity is measurable, but it is not a clean analogue of recovery because both dense runs have entered a degraded regime under AdamW at $10^{-4}$. A smaller learning rate might avoid this regime, but that hypothesis was not tested.

A public dense text-fine-tuned checkpoint trained with a different recipe provides a separate reference. Its gain is $-0.022$ \ci{-0.061}{+0.017} at $3.84$\,s and $+0.091$ \ci{+0.042}{+0.141} at $10.24$\,s, giving $J=+0.113$ \ci{+0.051}{+0.173}. Because the recipe and training history are not matched, this result is supporting context rather than a causal control.

\begin{table*}[t]
\centering
\caption{Recovery across duration and domain. All intervals are $95\%$ prompt-bootstrap intervals. Positive KL recovery means that \PFT{} reduces KL to the matched real reference relative to P. Human-CLAP, KL and PANNs sweep rescoring are corroborative analyses.}
\label{tab:evidence}
\setlength{\tabcolsep}{4.0pt}
\begin{tabular}{@{}lcccc@{}}
\toprule
Duration & CLAP $R$ & Human-CLAP $R$ & KL recovery & PANNs capture $R$ \\
\midrule
$3.84$ s & $.085$ $[.066,.105]$ & $.189$ $[.162,.217]$ & $.66$ $[.43,.92]$ & $.19$ $[.06,.32]$ \\
$5.12$ s & $.139$ $[.115,.164]$ & $.278$ $[.243,.314]$ & $1.16$ $[.91,1.41]$ & $.38$ $[.23,.53]$ \\
$7.68$ s & $.201$ $[.175,.227]$ & $.371$ $[.336,.405]$ & $1.95$ $[1.68,2.23]$ & $.76$ $[.61,.91]$ \\
$10.24$ s & $.244$ $[.215,.273]$ & $.375$ $[.340,.409]$ & $2.22$ $[1.92,2.52]$ & $.86$ $[.70,1.02]$ \\
$J$ & $.159$ $[.131,.187]$ & $.185$ $[.151,.220]$ & $1.56$ $[1.20,1.92]$ & $.67$ $[.49,.85]$ \\
\midrule
Domain & $n$ & $R(3.84\,\mathrm{s})$ & $R(10.24\,\mathrm{s})$ & $\rho_{\rm dense}(10.24\,\mathrm{s})$ \\
\midrule
AudioCaps & 192 & $.085$ $[.066,.105]$ & $.244$ $[.215,.273]$ & $.82$ \\
Clotho & 96 & $.098$ $[.072,.125]$ & $.210$ $[.176,.243]$ & $.74$ \\
Hip-hop & 127 & $.026$ $[.007,.044]$ & $.027$ $[.003,.050]$ & $.119$ $[.015,.215]$ \\
\bottomrule
\end{tabular}
\end{table*}

\subsection{Duration dependence weakens where recovery is weak}
Clotho shows that leaving AudioCaps does not by itself remove recovery. At $10.24$\,s, Clotho gives $R=+0.210$ \ci{+0.176}{+0.243} and closes $74\%$ of the dense gap. Its duration interaction remains substantial, $J=+0.112$ \ci{+0.079}{+0.146}. The difference from the AudioCaps interaction is $+0.047$ \ci{+0.003}{+0.090}.

Hip-hop behaves differently. With all $127$ prompts, recovery stays near $+0.03$ at both durations and its interaction is $J=+0.001$ \ci{-0.026}{+0.028}. Dense alignment remains about $+0.11$ above shuffled-caption chance, so the small gain is not caused by a floor-limited battery. Recovery closes about $12\%$ of the native dense gap, but this ratio is imprecise because both its small numerator and estimated denominator contribute uncertainty. Across the three domains, the duration interaction is largest where recovery itself is largest, intermediate on Clotho, and unresolved near zero on hip-hop. This connects the duration and domain effects rather than treating them as independent phenomena.

\subsection{Prompt-sample and short-generation diagnostics}
Severity~1 shows genuine prompt-set heterogeneity. The original outcome-blind $80$-prompt draw gives $J=+0.044$ \ci{-0.000}{+0.088}, whereas a disjoint $96$-prompt draw gives $+0.169$ \ci{+0.115}{+0.222}; their difference is $+0.124$ \ci{+0.058}{+0.194}. Both draws were selected without outcomes, but the second used a different seeded-hash salt and required five caption rows per clip, so they are not identical sampling rules. Severity~2 is much more stable under the same diagnostic. Its two outcome-blind $96$-prompt halves give $J=+0.159$ \ci{+0.116}{+0.200} and $+0.160$ \ci{+0.122}{+0.198}, with a difference of $+0.001$ \ci{-0.056}{+0.057}. Clotho supplies an additional independent prompt set on which the interaction also resolves.

Under CLAP, base-model degradation at $3.84$\,s does not account for the interaction. The dense model's floor-corrected response from $3.84$ to $10.24$\,s is $+0.142$ \ci{+0.111}{+0.172}, close to the real-audio crop response of $+0.150$ \ci{+0.133}{+0.166}. This is scorer-internal evidence rather than a perceptual claim. A separate crop analysis shows that recovery itself still depends on how the clip is generated. Scoring the first $3.84$\,s of a $10.24$\,s generation gives $R_{\rm crop}=+0.172$ \ci{+0.150}{+0.194}, which is $+0.087$ \ci{+0.065}{+0.110} above recovery from a separately generated $3.84$\,s clip.
